\documentclass[11pt]{article}

\usepackage{amsmath,amssymb}
\usepackage{xurl}
\usepackage[hidelinks]{hyperref}
\setlength{\emergencystretch}{2em}

% Shared commands for notes in docs/paper-gaps/.
% Notation follows the LDT paper (references/ldt-paper/).

\newcommand{\Lean}{\textsc{Lean}}
\newcommand{\leanid}[1]{\textsf{#1}}
\newcommand{\ghissue}[1]{\href{https://github.com/LionSR/MIPStarRE/issues/#1}{GitHub issue \##1}}
\newcommand{\ghpr}[1]{\href{https://github.com/LionSR/MIPStarRE/pull/#1}{GitHub PR \##1}}

% --- Paper-compatible notation ---
\newcommand{\F}{\mathbb{F}}
\newcommand{\N}{\mathbb{N}}
\newcommand{\C}{\mathbb{C}}
\newcommand{\tr}{\operatorname{tr}}
\newcommand{\ot}{\otimes}
\newcommand{\E}{\mathop{\bf E\/}}
\newcommand{\bx}{\mathbf{x}}
\newcommand{\by}{\mathbf{y}}
\newcommand{\bu}{\mathbf{u}}
\newcommand{\bv}{\mathbf{v}}

% Approximation relation (paper uses \approx_\eps and \simeq_\zeta)
\newcommand{\approxeps}{\approx_{\varepsilon}}
\newcommand{\simgeq}{\simeq_{\zeta}}

% Polynomial spaces (paper uses \polyfunc, \polysub, \polymeas)
\newcommand{\polyfunc}[3]{\operatorname{Poly}(#1,#2,#3)}
\newcommand{\polysub}[3]{\operatorname{PolySub}(#1,#2,#3)}
\newcommand{\polymeas}[3]{\operatorname{PolyMeas}(#1,#2,#3)}


\title{The Support of the Distinct-Tuple Distribution}
\author{}
\date{2026-05-01}

\begin{document}
\maketitle

\section{The assertion}

This note concerns the distinct-tuple distribution used in the pasting section of
\cite{Ji2020LowIndividualDegree}.%
\footnote{Tracked as part of the discrepancy audit in \ghissue{930}.  In the
TeX source consulted here, the definition and comparison proposition are in
\path{references/ldt-paper/ld-pasting.tex}, lines~167--213.}
For $k\ge 1$, the paper defines
\[
  \mathsf{Distinct}_k=\{(x_1,\ldots,x_k)\in \F_q^k: x_i\ne x_j
    \text{ for }i\ne j\}
\]
and writes $(\mathbf y_1,\ldots,\mathbf y_k)\sim \mathsf{Distinct}_k$ for a uniformly
random element of this set.  The following comparison proposition then bounds
the total variation distance between a uniformly random tuple in $\F_q^k$ and a
uniformly random tuple in $\mathsf{Distinct}_k$ by $k^2/q$.

\section{The point at issue}

The distribution on $\mathsf{Distinct}_k$ is a genuine probability distribution only
when $\mathsf{Distinct}_k$ is nonempty.  This requires $k\le q$.  The statement of the
comparison proposition and the surrounding pasting theorem do not state this
condition.  Later in the section the paper assumes $k\ge 400md$, but it does
not separately assume $k\le q$.

Thus, read literally, the phrase ``uniformly random element of
$\mathsf{Distinct}_k$'' is undefined in the case $k>q$.  In that case no injective map
from $\{1,\ldots,k\}$ to $\F_q$ exists.  The displayed bound $k^2/q$ is then
already at least $1$, so the intended comparison is only used in a regime where
the estimate is trivial, but the probability distribution itself still needs a
formal interpretation.

\section{The blueprint statement}

The blueprint repeats the paper's convention.  The definition
\path{def:distinct-tuples} in
\path{blueprint/src/chapter/ch09_pasting.tex}, lines~198--202, defines
$\mathsf{Distinct}_k$ for $k\ge 1$ and links it to the formal distinct-tuple
definitions.  The proposition \path{prop:ld-dnoteq} in the same file,
lines~204--226, then states the total-variation comparison between a uniformly
random tuple in $\F_q^k$ and a uniformly random tuple in
$\mathsf{Distinct}_k$.

As in the paper, the blueprint does not state the side condition $k\le q$.
Both blueprint nodes are marked as formalized, so the formalization supplies the
missing boundary convention: the Lean declaration linked from the blueprint is a
finite-support weight for all $k$, and it becomes a probability distribution
only in the nonempty-support case $k\le q$.

\section{The formal repair}

The formal development uses a finite-support weighted distribution rather than
forcing every distribution object to be normalized.  The declaration
\path{MIPStarRE.LDT.Pasting.distinctTupleDistribution} assigns weight
$1/|\mathsf{Distinct}_k|$ to each distinct tuple when the support is nonempty, and
weight zero outside the support.  If $k>q$, the support is empty, so this is the
zero finite measure rather than a probability distribution.%
\footnote{See \doclink{MIPStarRE/LDT/Pasting/Defs/Tuples.lean}, lines~23--36.}

The formalization records both facts.  The general theorem
\path{MIPStarRE.LDT.Pasting.distinctTupleDistribution_weight_sum_le_one}
proves that the total mass is at most $1$ for every $k$.  The normalized form
\path{MIPStarRE.LDT.Pasting.distinctTupleDistribution_weight_sum_eq_one_of_le}
adds the missing hypothesis $k\le q$ and proves that the total mass is exactly
$1$.%
\footnote{The normalized statement is in
\doclink{MIPStarRE/LDT/Pasting/BridgeLemmas/Common.lean}, lines~311--343 of the
source file.}

The total-variation comparison
\path{MIPStarRE.LDT.Pasting.ldDnoteq} is correspondingly case-split.  When
$k\le q$, it proves the usual collision-probability estimate.  When $k>q$, the
support is empty; the formal total variation between the uniform tuple
distribution and the zero distinct-tuple weight is $1/2$, and this is bounded
by $k^2/q$ because $k>q$.%
\footnote{See \doclink{MIPStarRE/LDT/Pasting/Core.lean}, lines~561--805.}
Downstream averaging lemmas use the same convention: they compare expectations
of nonnegative functions bounded by $1$, and the $k>q$ branch is discharged by
the now-trivial $k^2/q$ error term.

\section{Conclusion}

The paper implicitly treats the distinct-tuple distribution as if it were
available for all $k\ge 1$.  A literal probability distribution exists only for
$k\le q$.  The formal development makes the boundary case total by interpreting
the distinct-tuple object as a finite nonnegative weight for all $k$, and as a
probability distribution precisely under the additional hypothesis $k\le q$.
This is a harmless formal repair: it preserves the paper argument in the
nonempty-support case and makes the empty-support case a trivial bound rather
than an undefined sampling operation.

\bibliographystyle{alpha}
\bibliography{references}

\end{document}
